Speech Support for Angelman Syndrome
This project helps individuals with Angelman syndrome practice speaking through adaptive voice-based therapy, improve communication over time, and deliver progress reports to neurologists.

Inspiration
I was inspired to build this project after shadowing a neurologist and meeting a patient with Angelman syndrome. During the visit, I saw how limited verbal communication affected daily interactions, despite the patient being very expressive and engaged in other ways.

What it does
This speech support device is an adaptive, voice-based speech therapy system designed to help people with Angelman Syndrome continue practicing speech and vocalization. The system guides patients through spoken therapy activities such as memory recall, word association, and also sound practice. It listens to user responses, provides verbal feedback, tracks performance over time, and generates structured progress reports that can be shared with neurologists/speech-language pathologists.

How I built it
The project was built using Python and runs entirely from my mac terminal. Text-to-speech is used to deliver spoken instructions and feedback, while microphone input and speech recognition are used to detect and evaluate vocal responses. Each therapy session is logged, analyzed, and summarized into a clinician-ready report, which can be automatically sent via email (using gmail API).

Challenges I ran into
One major challenge was designing a system that encourages speaking without relying on typing or complex interfaces, since many individuals with Angelman syndrome are minimally verbal.

Accomplishments that I'm proud of
I'm proud of building a fully voice-driven therapy workflow that adapts in real time and produces meaningful, structured progress reports.

What I learned
I learned that effective speech therapy tools do not need to be too complicated to be impactful. Measuring engagement and improvement over time can be just as good as measuring accuracy.

What's next for Speech Support for Angelman Syndrome
Future plans include expanding therapy activities, improving speech recognition robustness, integrating alternative communication methods, and developing a caregiver-friendly interface.
